\documentclass[letterpaper,10pt]{article}
\usepackage[margin=0.3in]{geometry}
\usepackage{multicol}
\usepackage{amsmath,amssymb}
\usepackage{enumitem}
\usepackage{titlesec}
\usepackage{fancyhdr}
\usepackage{graphicx}

% Set font to sans-serif for clarity
\renewcommand{\familydefault}{\sfdefault}
\usepackage{helvet}

% Set section fonts to 7pt
\usepackage{anyfontsize}
\usepackage{sectsty}
\sectionfont{\fontsize{7}{8}\selectfont}
\subsectionfont{\fontsize{7}{8}\selectfont}
\subsubsectionfont{\fontsize{7}{8}\selectfont}

% Remove page numbers
\pagestyle{fancy}
\fancyhf{}
\renewcommand{\headrulewidth}{0pt}
\renewcommand{\footrulewidth}{0pt}

% Reduce space between sections
\titlespacing{\section}{0pt}{*0}{*0}
\titlespacing{\subsection}{0pt}{*0}{*0}
\titlespacing{\subsubsection}{0pt}{*0}{*0}

% Begin document
\begin{document}
\small
\begin{multicols}{3}

% P1: Key Statements for True/False Assessment
\section*{P1: Key Statements for T/F Assessment}

\textbf{Guidelines}: Use the following key statements and concepts to assess the correctness of true/false questions related to control systems.

\begin{itemize}[leftmargin=*]
    \item \textbf{Energy Storage vs. Dissipation}: A damper dissipates energy; a spring stores energy.
    
    \item \textbf{Damping Ratio (\(\zeta\))}: Underdamped systems have \(0 < \zeta < 1\); critically damped systems have \(\zeta = 1\); overdamped systems have \(\zeta > 1\).
    
    \item \textbf{Stability Criteria}: Systems are stable if all poles of the transfer function have negative real parts; poles in the right half-plane indicate instability.
    
    \item \textbf{System Responses}: A stable step response does not guarantee a stable ramp response; overshoot occurs only in underdamped systems (\(\zeta < 1\)); settling time (\(T_s\)) is typically greater than or equal to rise time (\(T_r\)).
    
    \item \textbf{Laplace Transform Properties}: Laplace transforms exist only for functions of exponential order; the Initial Value Theorem relates to \(t \to 0\), not \(t \to \infty\).
    
    \item \textbf{Response Types}: Impulse response is the output when an impulse is applied; system type is determined by the number of integrators (poles at the origin) in the open-loop transfer function; Type 0 systems have a non-zero steady-state error for step inputs; adding a zero can shape and potentially improve transient response; increasing feedback gain does not always improve stability; observability and controllability are dual concepts.
    
    \item \textbf{Frequency Response Analysis}: Phase margin is a frequency domain measure of stability, not a time-domain measure; Bode plots represent frequency-domain response, not time-domain; bandwidth is inversely proportional to rise time, meaning higher bandwidth typically implies a faster response; feedback control can stabilize certain unstable systems through proper design.
    
    \item \textbf{Routh-Hurwitz Criterion}: Determines the number of poles in the right half-plane but does not provide their exact locations.
\end{itemize}

% P2: Performance Indicators
\section*{P2: Performance Indicators}

\subsection*{First-Order Indicators}
\[
G(s) = \frac{K}{\tau s + 1}
\]
\begin{itemize}[leftmargin=*]
    \item \textbf{Rise Time (\(T_r\))}:
    \begin{itemize}[leftmargin=*]
            \item \( T_r \approx 2.2\tau \)
            \item Exact: \( y(T_r) = 0.9K \).
    \end{itemize}
    
    \item \textbf{Settling Time (\(T_s\))}:
    \begin{itemize}[leftmargin=*]
            \item \( T_s \approx \frac{4}{\tau} \) (2\% criterion)
            \item \( T_s \approx \frac{3}{\tau} \) (5\% criterion)
            \item Exact:  \( |y(T_s) - K| \leq 0.02K \).
    \end{itemize}
    
    \item \textbf{Steady-State Error (\(e_{ss}\))}:
        \begin{itemize}[leftmargin=*]
            \item For step input: \( e_{ss} = \frac{1}{1+K} \)
            \item \( e_{ss} = \lim_{s \to 0} s \cdot \frac{1}{1+G(s)} \cdot \frac{1}{s} = \frac{1}{1+K} \)
        \end{itemize}
\end{itemize}

\subsection*{Second-Order Indicators}
\[
G(s) = \frac{\omega_n^2}{s^2 + 2\zeta\omega_n s + \omega_n^2}
\]
\begin{itemize}[leftmargin=*]
    \item \textbf{Underdamped (\(0 < \zeta <1\))}:
    \begin{itemize}[leftmargin=*]
        \item \textbf{Damped Natural Frequency (\(\omega_d\))}: 
        \begin{itemize}[leftmargin=*]
            \item \( \omega_d = \omega_n\sqrt{1-\zeta^2} \)
        \end{itemize}
        
        \item \textbf{Peak Time (\(T_p\))}:
        \begin{itemize}[leftmargin=*]
                \item \( T_p = \frac{\pi}{\omega_d} \)
        \end{itemize}
        
        \item \textbf{Maximum Overshoot (\(M_p\))}:
        \begin{itemize}[leftmargin=*]
                \item \( M_p = e^{-\zeta\pi/\sqrt{1-\zeta^2}} \times 100\% \)
        \end{itemize}
        
        \item \textbf{Rise Time (\(T_r\))}:
        \begin{itemize}[leftmargin=*]
                \item \( T_r \approx \frac{\pi - \arccos(\zeta)}{\omega_d} \)
                \item Exact:  \( y(T_r) = 0.1K \) to \( y(T_r) = 0.9K \)
        \end{itemize}
        
        \item \textbf{Settling Time (\(T_s\))}:
        \begin{itemize}[leftmargin=*]
                \item \( T_s = \frac{4}{\zeta\omega_n} \) (2\% criterion)
                \item \( T_s = \frac{3}{\zeta\omega_n} \) (5\% criterion)
                \item Exact: \( |y(T_s) - K| \leq 0.02K \)
        \end{itemize}
    \end{itemize}
    
    \item \textbf{Critically Damped (\(\zeta =1\))}:
    \begin{itemize}[leftmargin=*]
        \item \textbf{Step Response}: 
        \begin{itemize}[leftmargin=*]
            \item \( y(t) = 1 - (1+\omega_n t)e^{-\omega_n t} \)
        \end{itemize}
        
        \item \textbf{Rise Time (\(T_r\))}:
        \begin{itemize}[leftmargin=*]
            \item \( T_r \approx \frac{1}{\omega_n} \)
            \item Exact: \( y(T_r) = 0.9K \)
        \end{itemize}
        
        \item \textbf{Settling Time (\(T_s\))}:
        \begin{itemize}[leftmargin=*]
            \item  \( T_s = \frac{4}{\omega_n} \)
        \end{itemize}
        
        \item \textbf{Maximum Overshoot (\(M_p\))}:
        \begin{itemize}[leftmargin=*]
            \item \( M_p = 0\% \). No overshoot
        \end{itemize}
    \end{itemize}
    
    \item \textbf{Overdamped (\(\zeta >1\))}:
    \begin{itemize}[leftmargin=*]
        \item \textbf{Rise Time (\(T_r\))}:
        \begin{itemize}[leftmargin=*]
            \item \( T_r \approx \frac{1.8}{\omega_n} \)
            \item Exact: \( y(T_r) = 0.1K \) to \( y(T_r) = 0.9K \)
        \end{itemize}
        
        \item \textbf{Settling Time (\(T_s\))}:
        \begin{itemize}[leftmargin=*]
            \item \( T_s \approx \frac{4}{\zeta\omega_n} \)
        \end{itemize}
        
        \item \textbf{Maximum Overshoot (\(M_p\))}:
        \begin{itemize}[leftmargin=*]
            \item  \( M_p = 0\% \). No overshoot
        \end{itemize}
    \end{itemize}
    
    \item \textbf{Resonant Frequency (\(\omega_r\))}:
    \begin{itemize}[leftmargin=*]
            \item \( \omega_r = \omega_n\sqrt{1-2\zeta^2} \)
    \end{itemize}
\end{itemize}


% P3: Closed-Loop Transfer Function
\section*{P3: Closed-Loop Transfer Function}
\[
T(s) = \frac{G(s)}{1 + G(s)H(s)}
\]
\textbf{Steps to Derive:}
\begin{enumerate}[leftmargin=*]
    \item Identify \( G(s) \) (open-loop transfer function) and \( H(s) \) (feedback transfer function).
    \item Apply the feedback formula: \( T(s) = \frac{G(s)}{1 + G(s)H(s)} \).
    \item Simplify the expression and convert to standard form if necessary.
\end{enumerate}

\noindent
\textbf{Block Diagram Reduction Techniques:}
\begin{itemize}[leftmargin=*]
    \item \textbf{Series Connection}: Multiply transfer functions.
    \item \textbf{Parallel Connection}: Add transfer functions.
    \item \textbf{Feedback Loop}: Use \( \frac{G}{1+GH} \).
    \item \textbf{Moving Blocks}: Rearrange blocks while maintaining system equivalence.
\end{itemize}

% P4: Root Locus
% P4: Root Locus
\section*{P4: Root Locus}

\subsection*{Construction Steps}

\begin{enumerate}[leftmargin=*]
    \item \textbf{Plot Poles and Zeros}: Mark all poles (\( \times \)) and zeros (\( \circ \)) of \( G(s)H(s) \) on the complex plane.
    
    \item \textbf{Determine Real Axis Segments}: Root locus exists on real axis segments with an odd number of poles and zeros to their right.
    
    \item \textbf{Calculate Asymptotes}:
    \begin{itemize}[leftmargin=*]
        \item \textbf{Number}: \( p - z \) (where \( p \) is number of poles and \( z \) is number of zeros).
        \item \textbf{Angles}: \( \theta = \frac{(2q + 1)180^\circ}{p - z} \), \( q = 0,1,\ldots,p-z-1 \).
        \item \textbf{Centroid}: \( \sigma = \frac{\sum \text{Poles} - \sum \text{Zeros}}{p - z} \).
    \end{itemize}
    
    \item \textbf{Identify Breakaway/Break-in Points}: Solve \( \frac{dK}{ds} = 0 \) and verify they lie on real axis segments.
    
    \item \textbf{Determine Departure/Arrival Angles}: Apply angle condition \( \sum \text{Angles from Poles} - \sum \text{Angles from Zeros} = 180^\circ + 360^\circ k \) for integer \( k \).
    
    \item \textbf{Sketch Root Locus}: Combine real segments, asymptotes, break points, and angle directions symmetrically.

    \item \textbf{Given a Pole Location} \( s = -\sigma \pm j\omega_d \):
    \[
    \omega_n = \sqrt{\sigma^2 + \omega_d^2}
    \]
    \[
    \zeta = \frac{\sigma}{\omega_n}
    \]
\end{enumerate}



% Bonus: PID
\section*{Bonus: PID Controllers}
\subsection*{PID Equation}
\[
G_c(s) = K_p + \frac{K_i}{s} + K_d s
\]
\subsection*{Effects of PID Components}
\begin{itemize}[leftmargin=*]
    \item \textbf{Proportional (P)}: Increases responsiveness; reduces steady-state error; excessive \( K_p \) can cause overshoot and potential instability.
    
    \item \textbf{Integral (I)}: Eliminates steady-state error; can reduce system stability; may cause increased overshoot.
    
    \item \textbf{Derivative (D)}: Improves system stability and damping; reduces overshoot; sensitive to measurement noise.
\end{itemize}

\subsection*{PID in Frequency Domain}
\[
G_c(s) = K_p \left(1 + \frac{1}{T_i s} + T_d s \right)
\]
\begin{itemize}[leftmargin=*]
    \item \( T_i = \frac{K_p}{K_i} \): Integral time constant.
    \item \( T_d = \frac{K_d}{K_p} \): Derivative time constant.
\end{itemize}

% Additional Content (Optional)
\section*{}
% You can add more sections or content here if needed

\end{multicols}
\end{document}
